% ============================================================
%  Entity Query Language: Syntax, Semantics, and Complexity
%  Standalone document — compile with pdflatex
% ============================================================

\documentclass[11pt, a4paper]{article}

% ---------- packages ----------
\usepackage[T1]{fontenc}
\usepackage[utf8]{inputenc}
\usepackage{geometry}
\geometry{margin=2.5cm}
\usepackage{amsmath, amssymb, amsthm}
\usepackage{mathtools}
% Semantic brackets, as in stmaryrd, without that package.
\newcommand{\llbracket}{\mathopen{[\![}}
\newcommand{\rrbracket}{\mathclose{]\!]}}
\usepackage{booktabs}
\usepackage{array}
\usepackage[dvipsnames]{xcolor}
\usepackage{listings}

% Python colour scheme
\definecolor{codebg}{RGB}{248,248,248}
\definecolor{kwcolor}{RGB}{0,100,180}
\definecolor{strcolor}{RGB}{60,140,60}
\definecolor{cmtcolor}{RGB}{130,130,130}
\definecolor{decocolor}{RGB}{160,60,160}
\definecolor{builtincolor}{RGB}{180,80,0}

\lstdefinestyle{eqlpython}{
  language=Python,
  backgroundcolor=\color{codebg},
  basicstyle=\small\ttfamily,
  keywordstyle=\color{kwcolor}\bfseries,
  keywordstyle=[2]\color{decocolor}\bfseries,
  keywords=[2]{dataclass},
  stringstyle=\color{strcolor},
  commentstyle=\color{cmtcolor}\itshape,
  emphstyle=\color{builtincolor},
  emph={True,False,None,List,any,sum,not,len},
  showstringspaces=false,
  breaklines=true,
  frame=none,
  xleftmargin=0pt,
  xrightmargin=0pt,
  aboveskip=6pt,
  belowskip=4pt,
}
\lstset{style=eqlpython}
\usepackage{hyperref}
\hypersetup{colorlinks=true, linkcolor=NavyBlue, citecolor=NavyBlue, urlcolor=NavyBlue}

% ---------- theorem environments ----------
\theoremstyle{plain}
\newtheorem{theorem}{Theorem}[section]
\newtheorem{proposition}[theorem]{Proposition}
\theoremstyle{definition}
\newtheorem{definition}[theorem]{Definition}
\theoremstyle{remark}
\newtheorem{remark}[theorem]{Remark}

% ---------- title ----------
\title{\textbf{Entity Query Language}\\[4pt]
       \large Syntax, Semantics, and Complexity}
\author{Anonymous Author(s)}
\date{\today}

% ============================================================
\begin{document}
% ============================================================

\maketitle
\tableofcontents
\bigskip

% ============================================================
\section{Overview}
\label{sec:eql:overview}
% ============================================================

Entity Query Language (EQL) is a symbolic, declarative, relational query language  with a First Order Logic like syntax that is embedded natively in
Python, designed to make Python object-oriented domain models directly queryable without any
additional representation layer. It separates symbolic, syntactical representation from evaluation semantics, thus allowing for different evaluation backends each with different semantics (e.g. open world vs closed world vs probabilistic inference). Furthermore, this layer separation increases maintainability of the language due to the separation of concerns between syntax and evaluation semantics.

The evaluation semantics of this document represent the default backend of EQL which obeys closed world reasoning that includes Domain Closure, and Negation as Failure. Future work will include other backends.

This document provides a formal specification of EQL's
syntax, semantics, and computational complexity.

\subsection{Motivation and Design Principles}

Existing query languages impose a representational gap between the domain model and the query
layer.  SQL operates over flat relational tables and requires either manual schema maintenance
or an Object-Relational Mapping (ORM) layer to bridge the gap to object-oriented code.
SPARQL operates over RDF triples, requiring a separate knowledge graph representation.  Datalog
and its extensions operate over symbolic facts that must be populated from the application
separately.  In all cases, the developer must maintain two representations of the same domain:
one in their programming language, one in the query system.

EQL eliminates this gap.  Its design rests on the following principles:

\begin{enumerate}
  \item \textbf{No Impedence Mismatch and No Symbol Grounding Problems.}  EQL operates directly on live Python objects.
        There is no schema to define, no tables to populate, and no ORM to configure.  The
        object-oriented domain model the developer already writes \emph{is} the knowledge base.

  \item \textbf{Executable queries as first-class Python.}  Queries are Python expressions, not
        strings passed to an external compiler.  This means queries are statically analysable
        by development tools, debuggable at runtime, and subject to normal Python error
        handling and type checking.  There is no compilation boundary between query and
        application code.

  \item \textbf{General programs as query components.}  Any Python callable---including
        arbitrary programs---can participate in a query as a predicate or symbolic function,
        subject to determinism and purity requirements.  This goes beyond the built-in
        predicates of SQL or the fixed vocabulary of SPARQL.

  \item \textbf{Structural depth.}  EQL queries traverse nested objects, lists, and
        dictionaries natively.  There is no restriction to flat relational structure.  An
        attribute of a field that is itself a list of objects is queried by index or method
        access, not by a join.

  \item \textbf{FOL axioms in class definitions.}  Predicates are Python classes, and their
        characteristic functions can encode arbitrary first-order-logic constraints over the
        domain.  These axioms live directly in the class definition alongside the data model,
        making the logical structure of the domain inspectable and revisable.
\end{enumerate}

\subsection{Relation to Existing Query Languages}

EQL's logical foundations are drawn from conjunctive query evaluation over finite relational
structures~\cite{chandra1977optimal,abiteboul1995foundations}, Negation as
Failure~\cite{clark1978negation}, and the TBox/ABox distinction of Description
Logics~\cite{baader2003description}.  It differs from these traditions in three key respects.

\emph{SQL and relational databases}: SQL requires a separate schema, table-based
storage, and an ORM layer to map objects to relations.  EQL queries the object model directly;
there are no tables and no joins -- object traversal replaces join navigation.

\emph{Datalog}: Datalog rules are stored in the knowledge base and applied to a
fixpoint.  In EQL, rules are user-invoked programs; there is no fixpoint.  EQL's fact base is always a flat finite set of ground object instances. Even though rules exist they are not automatically invoked when answering queries unless they are called as part of the query.

\emph{SPARQL and RDF}: SPARQL requires a triple-store representation distinct from the
application model.  EQL operates on the application's own Python objects with no translation
step.

\paragraph{Contributions.}
This document makes the following formal contributions:
\begin{enumerate}
  \item A typed signature $\Sigma = (\mathcal{S}, \mathcal{C}, \mathcal{M}, \mathcal{P},
        \mathcal{F})$ for EQL with a formal TBox/ABox knowledge base structure and a
        subsumption ordering on class symbols.
  \item A distinction between query variables $\mathcal{V}_Q$ (user-supplied finite domains,
        optional class annotation) and inference variables $\mathcal{V}_I$ (class-annotated,
        domain populated from rule body answer sets), formalising EQL's rule mechanism.
  \item An abstract syntax covering terms with structural depth (field, index, and method
        access), three atom forms including Boolean term atoms $[t]_\mathbb{B}$, body formulas,
        conjunctive queries, unions, result quantifiers, and rules.
  \item A denotational semantics grounded in Clark's completion for negation as failure and
        the $\gamma$ operator of relational algebra for aggregation.
  \item Computational complexity bounds for progressively richer EQL query classes,
        distinguishing data complexity from combined complexity.
\end{enumerate}
% ============================================================
\section{Vocabulary and the Knowledge Base}
\label{sec:eql:vocab}
% ============================================================

\subsection{Signature}

A \emph{signature} $\Sigma = (\mathcal{S}, \mathcal{C}, \mathcal{M}, \mathcal{P}, \mathcal{F})$
consists of non-logical symbols over which the EQL language is built.  Variable symbols
$\mathcal{V}$ are logical symbols introduced separately in Section~\ref{sec:eql:vars}.

\paragraph{Class symbols.}
A finite set of \emph{class symbols} $\mathcal{S} = \{S_1, S_2, \ldots, S_k\}$ corresponds to
Python classes.  The class hierarchy defined by the developer induces a \emph{subsumption
ordering} $\sqsubseteq$ on $\mathcal{S}$, where $S' \sqsubseteq S$ if and only if $S'$ is a
subclass of $S$ in Python's class system (i.e., \texttt{issubclass($S'$, $S$)} holds).

\paragraph{Constant symbols.}
A finite set of \emph{constant symbols} $\mathcal{C}$ denotes the Python object instances
populating the domain.  Each constant is typed by one or more classes via the class hierarchy.

\paragraph{Mapped symbols.}
A finite set of \emph{mapped symbols} $\mathcal{M}$.  Each $m \in \mathcal{M}$ is a named
attribute of a class $S \in \mathcal{S}$, representing a property that is intrinsic to $S$
and accessible on any instance of $S$ via dot notation. They are referred to as Mapped symbols because they can be seen as a mapping applied to S. The attribute $m$ may be:

\begin{itemize}
  \item a \emph{field}, typed $m : S \to T$, returning a value of type $T$.
  \item an \emph{indexable field}, typed $m : S \to \mathrm{Indexable}[K, T]$, supporting
        element access by key $k \in K$, where $K$ is any hashable key type (e.g.,
        $\mathbb{Z}$ (i.e., integers) for lists and tuples, or an arbitrary hashable type for dictionaries).
  \item a \emph{method}, typed $m : S \times T_1 \times \cdots \times T_n \to T$, invocable
        with arguments.
\end{itemize}

\noindent
How these attributes are accessed syntactically is defined in the term grammar (Section~\ref{sec:eql:syntax}).

\paragraph{Predicate symbols.}
A finite set of \emph{predicate class symbols} $\mathcal{P}$.  Each $p \in \mathcal{P}$ is a
Python class that is a subclass of \texttt{Predicate} and implements a characteristic function
via \texttt{\_\_call\_\_(self) -> bool}.  Predicates occupy two roles:

\begin{enumerate}
  \item \emph{Constraint role} (query body): the expression $p(\bar{t})$ is an atomic formula
        whose truth value is determined at evaluation time by Boolean-casting the ground
        predicate instance (Section~\ref{sec:eql:atoms_eval}).

  \item \emph{Domain role} (fact base): instances of $p$ are ordinary Python objects and may be
        stored in $\mathcal{D}$, queried over, and used as arguments in other atoms.
\end{enumerate}

\paragraph{Function symbols.}
A finite set of \emph{symbolic functions} $\mathcal{F}$.  Each $f \in \mathcal{F}$ is a
Python callable of one of the following forms:
\begin{enumerate}
    \item A Python function decorated with \texttt{@symbolic\_function}.
    \item A class method accessed through an EQL variable.
\end{enumerate}
\noindent
Function outputs are not automatically stored in $\mathcal{D}$.
\newline

Every $f \in \mathcal{F}$ and every characteristic function $p \in \mathcal{P}$ is assumed
\textbf{pure} (no side effects on $\mathcal{D}$ or shared state) and \textbf{deterministic}
(same inputs always produce the same output), and \textbf{terminating} on every input
that evaluation passes to it.  These conditions are required at all times to ensure
correctness and termination of query evaluation.  An acyclic call graph over
$\mathcal{F} \cup \mathcal{P}$ rules out unbounded mutual recursion between callables, but it
does not imply termination, since a single non-recursive callable may loop.  These are
semantic preconditions (Section~\ref{sec:eql:eval}); they are not statically checkable for
arbitrary Python callables~\cite{rice1953classes}.

% ------------------------------------------------------------
\subsection{The Knowledge Base}
% ------------------------------------------------------------

\begin{definition}[EQL Knowledge Base]
\label{def:kb}
An EQL \emph{knowledge base} is a pair $\mathcal{K} = (\mathcal{T},\,\mathcal{D})$ where:
\begin{itemize}
  \item $\mathcal{T}$ \emph{(TBox)} is the Python class hierarchy: the set of class symbols
        $\mathcal{S}$, their attribute symbols $\mathcal{M}$, and the subsumption ordering
        $\sqsubseteq$.

  \item $\mathcal{D}$ \emph{(ABox)} is a finite set of ground typed object
        instances.  $\mathcal{D}$ contains no disjunctive
        assertions, no conditional facts, and no evaluatable statements.
\end{itemize}
\end{definition}

% ------------------------------------------------------------
\subsection{Variable Symbols}
\label{sec:eql:vars}
% ------------------------------------------------------------

Let $\mathcal{V} = \mathcal{V}_Q \cup \mathcal{V}_I$ be the set of \emph{variable symbols}. These are the logical symbols of EQL, ranging over domains of constants drawn from $\mathcal{C}$.
EQL provides two types of variables:

\emph{Query variables} $\mathcal{V}_Q$: each carries a finite \emph{domain} $D \subseteq
\mathcal{C}$, which is the only mandatory annotation.  A class annotation $S \in \mathcal{S}$
is optional; when supplied it constrains which objects are admissible in $D$. The domain
can be user-supplied or if $S$ is given and no $D$ is provided, then the full class extension $\mathrm{dom}(S)$ is used as $D$ instead.
Query variables are used in query bodies and heads.

\emph{Inference variables} $\mathcal{V}_I$: each carries a class annotation $S \in
\mathcal{S}$ and
constructor keyword arguments $\bar{k}$ mapping field names to body variables.  An inference
variable carries no pre-specified domain; its domain is populated at rule invocation time:
\[
  D_{x_I} \;=\;
  \bigl\{\, H(\bar{k}(\bar{a})) \;\bigm|\; \bar{a} \in \mathrm{ans}(Q_{\mathrm{body}},\,\mathcal{K}) \,\bigr\}
\]
In EQL, inference variables are constructed via \texttt{inference(H)(**kwargs)}.

% ============================================================
\section{Syntax}
\label{sec:eql:syntax}
% ============================================================

This section defines the \emph{abstract syntax} of EQL: the grammar of well-formed terms,
atoms, formulas, queries, and rules, independently of their Python surface representation.
The Python API constitutes a concrete realisation of this abstract language; where relevant,
correspondences to specific API constructs are noted as glosses (e.g., $\wedge$ is expressed
in EQL via \texttt{and\_()} or multiple arguments to \texttt{.where()}).

\subsection{Terms}

Terms are defined inductively over $\Sigma$:
\[
  t \;::=\;
    c
  \;\mid\; x
  \;\mid\; t.m
  \;\mid\; t.m[i]
  \;\mid\; t.m(\bar{t})
  \;\mid\; f(t_1, \ldots, t_n)
\]
where:
\begin{itemize}
  \item $c \in \mathcal{C}$ is a constant (ground object), also called a \emph{Literal} in
        EQL; e.g., a specific \texttt{Robot} instance.
  \item $x \in \mathcal{V}$ is a domain-annotated and (optionally) class-annotated variable;
        e.g., \texttt{r = variable(Robot, domain=robots)}.
  \item $t.m$ is a field access on term $t$, returning the value of field $m$; e.g.,
        \texttt{r.battery}.  Access chains such as $t.m.b$ are admitted, since $t.m$ is
        itself a term; e.g., \texttt{r.tasks[0].completed}.
  \item $t.m[i]$ is an index access retrieving the element at position $i$ (or with key $i$
        for a dictionary) of the indexable field $m$ of $t$; e.g., \texttt{r.tasks[0]}.
  \item $t.m(\bar{t})$ is a method call invoking method $m$ on $t$ with arguments $\bar{t}$;
        when $t$ or any element of $\bar{t}$ is an unground variable, the call is deferred;
        e.g., \texttt{r.name.startswith("R")}.
  \item $f(t_1, \ldots, t_n)$ is an explicit symbolic function application; when any $t_i$ is
        an unground variable, evaluation is deferred; e.g., \texttt{distance(r.position, base)}.
\end{itemize}

\subsection{Atomic Formulas}
\label{sec:eql:atoms}

An \emph{atom} has one of the following forms:
\[
  \alpha \;::=\;
    p(t_1, \ldots, t_n)
  \;\mid\;
    t_1 \bowtie t_2
  \;\mid\;
    [t]_\mathbb{B}
\]

\begin{itemize}
  \item $p(t_1,\ldots,t_n)$ is a \emph{predicate atom}: it denotes the instance of predicate
        class $p$ constructed with terms $t_1,\ldots,t_n$ as arguments.  Its truth value is
        determined at evaluation time by Boolean-casting the ground instance
        (Section~\ref{sec:eql:atoms_eval}); e.g., \texttt{HasIncompleteTask(r)}.

  \item $t_1 \bowtie t_2$ where
        ${\bowtie}\in\{=,\,\neq,\,<,\,\leq,\,>,\,\geq,\,\in\}$ is a \emph{comparison atom}
        over two terms; e.g., \texttt{r.battery > 50}.

  \item $[t]_\mathbb{B}$ is a \emph{Boolean term atom}: any term $t$ that evaluates to a
        Boolean value may appear as an atom by Boolean-casting its value.  This covers symbolic
        function calls ($f(t_1,\ldots,t_n)$ where $f$ returns \texttt{bool}), method calls
        ($t.m(\bar{u})$ returning \texttt{bool}), and Boolean field accesses ($t.m$ of Boolean
        type); e.g., \texttt{r.name.startswith("R")} or \texttt{is\_valid(r)}.
\end{itemize}

\begin{remark}[Built-in predicates]
EQL provides built-in predicate implementations for common TBox queries, including
\texttt{HasType}, which wraps \texttt{isinstance}, and \texttt{IsSubClass}, which wraps
\texttt{issubclass}.  These are members of $\mathcal{P}$ like any user-defined predicate and
receive no special syntactic treatment.
\end{remark}

\subsection{Formulas (Query Bodies)}

A \emph{body formula} $\phi$ is built from atoms by:
\[
  \phi \;::=\;
    \alpha
  \;\mid\; \phi_1 \wedge \phi_2
  \;\mid\; \phi_1 \vee \phi_2
  \;\mid\; \neg\phi
  \;\mid\; \exists x.\,\phi
  \;\mid\; \forall x.\,\phi
\]

\begin{itemize}
  \item Conjunctions ($\wedge$) are expressed via multiple arguments to \texttt{.where()} or
        \texttt{and\_()}.
  \item Disjunctions ($\vee$) are expressed via \texttt{or\_()}.
  \item Negation ($\neg$) is interpreted as Negation as Failure
        (Section~\ref{sec:eql:naf}), expressed via \texttt{not\_()}.
  \item Existential quantification ($\exists$) arises implicitly: a domain-annotated variable
        $x$ with domain $D$ acts as $\exists x \in D$ in the body.
  \item Universal quantification ($\forall$) is explicit and restricted to the body; it is
        evaluated over the finite closed domain $D$.
  \item The domain of a quantified variable may depend on outer variables, for example the courses of a student $s$, written \texttt{flat\_variable(s.takes\_course)}.  Such a quantifier is evaluated
        separately for each binding of the outer variables; the outer variables are free in
        the quantified formula, not quantified by it.
\end{itemize}

\subsection{Queries}

A \emph{conjunctive query} (CQ) is:
\[
  Q(\bar{x}) \;\leftarrow\; \phi(\bar{x},\,\bar{y})
\]
where $\bar{x}$ are the head (output) variables, $\bar{y}$ are body-only existentially
quantified variables, and $\phi$ is a body formula.

A \emph{union of conjunctive queries} (UCQ) is a finite disjunction of CQs sharing the same
head variables, expressed in EQL via \texttt{or\_()} in a single \texttt{.where()} call:
\[
  Q(\bar{x}) \;\leftarrow\;
    \phi_1(\bar{x}) \vee \cdots \vee \phi_k(\bar{x})
\]
The same-head requirement ensures the result is a well-typed relation: all disjuncts
contribute tuples of the same arity and type.  In EQL this constraint is enforced
architecturally: \texttt{or\_()} operates inside the body formula $\phi$ of a single
\texttt{entity()} or \texttt{set\_of()} call, so the head variables are shared by
construction across all disjuncts.

\subsection{Result Quantifiers}
\label{sec:eql:quantifiers}

A \emph{result quantifier} wraps a query and enforces a cardinality constraint on its answer
set, producing a \emph{quantified query} that is itself a query.  The answer set is computed
normally; the constraint is then checked against the cardinality of that answer set.  If the
constraint is violated, a \texttt{QuantificationError} is raised rather than returning results.

A \emph{cardinality constraint} $C : \mathbb{N} \to \{\top, \bot\}$ is one of:
\[
  C \;::=\;
    \texttt{AtLeast}(n)
  \;\mid\; \texttt{AtMost}(n)
  \;\mid\; \texttt{Exactly}(n)
  \;\mid\; \texttt{Range}(m, n)
\]
with the obvious interpretations over the cardinality $|\mathrm{ans}(Q, \mathcal{K})|$.

A \emph{quantified query} pairs a query $Q$ with a constraint $C$:
\[
  \mathcal{Q}(Q,\, C) \;=\;
  \begin{cases}
    \mathrm{ans}(Q,\,\mathcal{K}) & \text{if } C\!\bigl(|\mathrm{ans}(Q,\,\mathcal{K})|\bigr) \\
    \texttt{QuantificationError} & \text{otherwise}
  \end{cases}
\]

EQL provides two quantifier constructors.  \texttt{an} is the general form: without a
constraint argument it applies the trivially satisfied constraint $C(k) = \top$ for all $k$,
acting as a declarative annotation that zero or more results are expected.  With a constraint
argument it enforces the given $C$.  \texttt{the} asserts that exactly one result exists and
is equivalent to \texttt{an} with $\texttt{Exactly}(1)$; it corresponds to Russell's definite
description $\iota x.\,\phi(x)$~\cite{whitehead1910principia}, which denotes the unique
individual satisfying $\phi$ and is undefined if no such individual exists or more than one
does.

\subsection{Rules}
\label{sec:eql:rules}

A \emph{rule} pairs a query body with an inference variable that constructs new instances from
the body's answer set:
\[
  x_I \;\leftarrow\; \phi(\bar{x},\,\bar{y})
  \qquad x_I \in \mathcal{V}_I,\quad \bar{x} \in \mathcal{V}_Q
\]
Rules are not stored in $\mathcal{K}$.  On user invocation, the engine evaluates $\phi$ over
the current $\mathcal{D}$, uses the answer set to populate $D_{x_I}$, and returns the
resulting instances.  The user may add them to $\mathcal{D}$.  There is no fixpoint iteration;
each invocation is a single query evaluation.

\section{Illustrative Example}
\label{sec:eql:example}
% ============================================================

The following example is given to ground the abstract definitions in Sections~\ref{sec:eql:syntax}
and~\ref{sec:eql:semantics}.  The domain in this example
consists of two classes:

\begin{lstlisting}
@dataclass
class Robot:
    name: str
    battery: int
    tasks: List[Task]

@dataclass
class Task:
    name: str
    completed: bool
\end{lstlisting}

\noindent
The fact base $\mathcal{D}$ contains three \texttt{Robot} instances and two \texttt{Task}
instances.  The class hierarchy $\mathcal{T}$ contains \texttt{Robot} and \texttt{Task} as
class symbols.

\paragraph{Example query.}
Retrieve all robots whose battery exceeds 50 and whose first task is not yet completed:

\begin{lstlisting}
r = variable(Robot, domain=robots)
results = an(entity(r).where(r.battery > 50, not_(r.tasks[0].completed))).tolist()
\end{lstlisting}

\noindent
In abstract syntax, this is the CQ $Q(r) \leftarrow \phi(r)$ where:
\[
  \phi(r) \;=\; (r.\mathtt{battery} > 50)
            \;\wedge\; \neg\,(r.\mathtt{tasks}[0].\mathtt{completed})
\]
Here $r \in \mathcal{V}_Q$, $r.\mathtt{battery}$ is a field access term
(Section~\ref{sec:eql:syntax}), $r.\mathtt{tasks}[0].\mathtt{completed}$ is an index access
chained with a field access, $r.\mathtt{battery} > 50$ is a comparison atom, and $\neg$ is
NAF negation (Section~\ref{sec:eql:naf}).

\paragraph{Example rule.}
Derive \texttt{HighLoadRobot} instances for robots whose battery is below 30 and that have at
least one incomplete task:

\begin{lstlisting}
@dataclass
class HighLoadRobot:
    robot: Robot

@dataclass(eq=False)
class HasIncompleteTask(Predicate):
    robot: Robot
    def __call__(self) -> bool:
        return any(not t.completed for t in self.robot.tasks)

r = variable(Robot, domain=robots)
high_load = inference(HighLoadRobot)(robot=r)
rule = an(entity(high_load).where(r.battery < 30, HasIncompleteTask(r)))
\end{lstlisting}

\noindent
In abstract syntax this is the rule $x_I \leftarrow \phi(r)$ where:
\[
  \phi(r) \;=\; (r.\mathtt{battery} < 30) \;\wedge\; \mathtt{HasIncompleteTask}(r)
\]
Here $r \in \mathcal{V}_Q$ is a query variable, $x_I \in \mathcal{V}_I$ is an inference
variable of class \texttt{HighLoadRobot} with argument mapping $\mathtt{robot} \mapsto r$
(Section~\ref{sec:eql:vars}), and $\mathtt{HasIncompleteTask}(r)$ is a predicate atom in
the body (Section~\ref{sec:eql:atoms}).  For each binding of $r$ satisfying $\phi$, a new
\texttt{HighLoadRobot} instance is constructed and forms the rule's answer set $D_{x_I}$.

% ============================================================

% ============================================================
\section{Semantics}
\label{sec:eql:semantics}
% ============================================================

\subsection{Domain Closure and Unique Names}

\begin{definition}[Domain Closure Assumption]
\label{def:dca}
Let $\tau : \mathcal{C} \to \mathcal{S}$ be a \emph{typing function} that assigns each
constant its most specific class.  The \emph{domain} of class $S \in \mathcal{S}$ is:
\[
  \mathrm{dom}(S) \;=\; \{\,c \in \mathcal{C} \mid \tau(c) \sqsubseteq S\,\}
\]
where $\sqsubseteq$ is the subsumption ordering of $\mathcal{T}$.  No individuals exist beyond
those in $\mathcal{D}$.
\end{definition}

\noindent
In EQL, $\tau$ is realised by Python's class system and domain membership $\tau(c) \sqsubseteq S$
can be checked at runtime using \texttt{isinstance}(\texttt{c},~\texttt{S}).

Every query variable $x \in \mathcal{V}_Q$ carries an explicit finite domain, enforcing the
DCA at query construction time.  Inference variables $x_I \in \mathcal{V}_I$ are domain-free
at construction; their domain is populated from the rule body's answer set at invocation time.
Distinct Python object identities denote distinct individuals (\emph{Unique Name Assumption}).
Together these fix a canonical finite interpretation of
$\mathcal{D}$~\cite{abiteboul1995foundations} over which evaluation proceeds.

\subsection{Evaluation Semantics}
\label{sec:eql:eval}

EQL realises \emph{conjunctive query evaluation by left-deep nested-loop
join}~\cite{chandra1977optimal,abiteboul1995foundations,vardi1982complexity}.  Because all
domains are finite, evaluation makes finitely many callable invocations
(Section~\ref{sec:eql:complexity}), so it terminates whenever every invoked callable
terminates (Section~\ref{sec:eql:vocab}).

Given $Q(\bar{x}) \leftarrow \phi(\bar{x}, \bar{y})$ and $\mathcal{K} = (\mathcal{T},
\mathcal{D})$, evaluation proceeds left to right through the expression tree of $\phi$:

\begin{enumerate}
  \item \textbf{Lazy variable binding.}  Variables are not enumerated upfront.  Each variable
        $x_i \in \mathcal{V}_Q$ is bound when its domain $D_i$ is first encountered in the
        left-to-right traversal of $\phi$.  The engine iterates over $D_i$ in a nested loop,
        passing each candidate binding forward to the next subexpression.

  \item \textbf{Short-circuit filtering.}  Under $\wedge$: if a subexpression fails for the
        current binding $\theta$, the remaining conjuncts are not evaluated for that $\theta$,
        and the next candidate is tried.  Under $\vee$: if the left disjunct succeeds, the
        right is skipped.  No enumeration is performed beyond what the expression requires.

  \item \textbf{Head projection.}  After $\phi$ is satisfied, any head variable $x_i \in
        \bar{x}$ not constrained by $\phi$ is joined with the current bindings via a Cartesian
        product over $D_i$ to produce complete result tuples.

  \item \textbf{Result collection.}  All bindings $\theta$ for which $\phi$ evaluates to true
        are collected; projecting onto $\bar{x}$ yields the answer set.
\end{enumerate}

\subsection{Answer Set}

\begin{definition}[Answer Set]
\label{def:ans}
The \emph{answer set} of query $Q(\bar{x}) \leftarrow \phi(\bar{x}, \bar{y})$ with respect to
$\mathcal{K}$ is:
\[
  \mathrm{ans}(Q,\,\mathcal{K})
  \;=\;
  \bigl\{\,\bar{a} \in D_{\bar{x}}
    \;\bigm|\;
    \mathcal{D} \models_{\mathrm{NAF}} \exists\bar{y}.\,\phi(\bar{a},\,\bar{y})
  \,\bigr\}
\]
where $D_{\bar{x}} = D_{x_1} \times \cdots \times D_{x_k}$ is the product of the
user-supplied domains of the variables in $\bar{x}$, and $\models_{\mathrm{NAF}}$ denotes
derivability under the Closed World Assumption with Negation as Failure
(Section~\ref{sec:eql:naf}).
\end{definition}

\subsection{Aggregation Semantics}
\label{sec:eql:aggregation}

An \emph{aggregated query} is a semantic transformation applied to the answer set of a CQ,
corresponding to the $\gamma$ operator of relational
algebra~\cite{klug1982equivalence,abiteboul1995foundations} and to SQL's \texttt{GROUP BY}
pattern.

Let $\mathcal{A}_f = \{\mathrm{count},\, \mathrm{sum},\, \mathrm{avg},\, \mathrm{max},\,
\mathrm{min},\, \mathrm{mode},\, \mathrm{multimode}\}$ be the set of aggregate functions,
each mapping a multiset of values to a scalar or set of scalars.  Given a CQ
$Q(\bar{x}) \leftarrow \phi(\bar{x}, \bar{y})$, a grouping key $\bar{g} \subseteq \bar{x}$,
and aggregate expressions $f_1(e_1), \ldots, f_k(e_k)$ where each $e_i$ is a term over
$\bar{x}$, the aggregated query evaluates as:
\[
  \gamma_{\bar{g};\; e_1 \to f_1,\,\ldots,\,e_k \to f_k}
  \!\bigl(\mathrm{ans}(Q,\mathcal{K})\bigr)
\]
This partitions $\mathrm{ans}(Q, \mathcal{K})$ by $\bar{g}$ and computes each $f_i$ over the
multiset of $e_i$ values within each group.  In EQL this is expressed via
\texttt{set\_of(...).grouped\_by($\bar{g}$)}.  Post-aggregation filtering via
\texttt{.having()} applies a condition after aggregation, equivalent to SQL's \texttt{HAVING};
this is distinct from \texttt{.where()}, which filters tuples before grouping as part of
$\phi$.  When no grouping key is specified, $f$ is computed over the entire answer set.

\subsection{Negation as Failure}
\label{sec:eql:naf}

EQL supports \emph{Negation as Failure} (NAF)~\cite{clark1978negation}.  The semantics of
$\models_{\mathrm{NAF}}$ used in Definition~\ref{def:ans} is given by Reiter's \emph{Closed
World Assumption} (CWA)~\cite{reiter1978closed}: let $\mathrm{CWA}(\mathcal{D})$ be the theory
obtained from the fact base $\mathcal{D}$ by adding $\neg\alpha$ for each ground atom
$\alpha$ that is not entailed by $\mathcal{D}$.  Then:
\[
  \mathcal{D} \models_{\mathrm{NAF}} \phi
  \;\iff\;
  \mathrm{CWA}(\mathcal{D}) \models \phi
\]
Because $\mathcal{D}$ contains only ground facts and no rules, $\mathrm{CWA}(\mathcal{D})$
has a single minimal model, and the CWA coincides with Clark's completion of $\mathcal{D}$.
Operationally, $\neg\phi$ succeeds in a query body if and only if $\phi$ fails finitely during
evaluation.  NAF is expressed in EQL via \texttt{not\_()}.

NAF is \textbf{non-monotonic}: adding new facts to $\mathcal{D}$ may invalidate previously
derived negative conclusions.  EQL does not provide monotonic or paraconsistent negation.

\subsection{Semantics of Terms}
\label{sec:eql:terms_eval}
Under a ground substitution $\theta$, terms evaluate as follows.

\paragraph{Field access.}
\[
  \llbracket t.m \rrbracket_\theta
  \;=\;
  \text{the value of field } m \text{ on object } \llbracket t \rrbracket_\theta
\]
Access chains evaluate left-associatively:
$\llbracket t.m.b \rrbracket_\theta = \llbracket (t.m).b \rrbracket_\theta$.

\paragraph{Index access.}
\[
  \llbracket t.m[i] \rrbracket_\theta
  \;=\;
  \llbracket t.m \rrbracket_\theta [\,\llbracket i \rrbracket_\theta\,]
\]
i.e., element $i$ of the indexable value of field $m$ on
$\llbracket t \rrbracket_\theta$.

\paragraph{Method call.}
\[
  \llbracket t.m(\bar{u}) \rrbracket_\theta
  \;=\;
  \llbracket t.m \rrbracket_\theta\!
  \bigl(\llbracket \bar{u} \rrbracket_\theta\bigr)
\]
i.e., method $m$ of $\llbracket t \rrbracket_\theta$ called with ground arguments
$\llbracket \bar{u} \rrbracket_\theta$.  If $t$ or any element of $\bar{u}$ is unground,
the call is deferred as a symbolic function application, returning a new EQL variable.

\paragraph{Symbolic function.}
\[
  \llbracket f(t_1, \ldots, t_n) \rrbracket_\theta
  \;=\;
  f_{\mathrm{impl}}\!\bigl(
    \llbracket t_1 \rrbracket_\theta,\,\ldots,\,
    \llbracket t_n \rrbracket_\theta
  \bigr)
\]
where $f_{\mathrm{impl}}$ is the underlying callable.  If any $t_i$ is unground, evaluation is
deferred.  Determinism and purity ensure the result is unique and $\mathcal{D}$ is unmodified.

% ------------------------------------------------------------
\subsection{Semantics of Atoms}
\label{sec:eql:atoms_eval}
% ------------------------------------------------------------

Atom evaluation determines whether a ground atom holds under substitution $\theta$.  Term
evaluation (Section~\ref{sec:eql:terms_eval}) is applied first to obtain ground values; truth determination is a separate
subsequent step.

\paragraph{Predicate atom.}
Applying $\theta$ to $p(t_1, \ldots, t_n)$ yields the ground predicate instance:
\[
  \llbracket p(t_1, \ldots, t_n) \rrbracket_\theta
  \;=\;
  p\!\bigl(
    \llbracket t_1 \rrbracket_\theta,\,\ldots,\,
    \llbracket t_n \rrbracket_\theta
  \bigr)
\]
The atom \emph{holds} if and only if the Boolean cast of this instance returns \texttt{True}:
\[
  \alpha \text{ holds under } \theta
  \;\iff\;
  \texttt{bool}\!\bigl(\llbracket p(t_1, \ldots, t_n) \rrbracket_\theta\bigr)
  = \texttt{True}
\]
The Boolean cast invokes the characteristic function \texttt{\_\_call\_\_} of the instance.

\paragraph{Comparison atom.}
The atom $t_1 \bowtie t_2$ holds under $\theta$ if and only if:
\[
  \llbracket t_1 \rrbracket_\theta \;\bowtie\; \llbracket t_2 \rrbracket_\theta
\]

\paragraph{Boolean term atom.}
The atom $[t]_\mathbb{B}$ holds under $\theta$ if and only if:
\[
  \texttt{bool}\!\bigl(\llbracket t \rrbracket_\theta\bigr) = \texttt{True}
\]
This applies to any term evaluating to a Boolean value: symbolic function calls, method calls,
and Boolean field accesses are all admitted as atoms via this form.

% ============================================================
\section{Computational Complexity}
\label{sec:eql:complexity}
% ============================================================

\emph{Data complexity} fixes the query and measures cost as $|\mathcal{D}|$ grows---the
relevant measure for a deployed system.  \emph{Combined complexity} lets both the query and
$\mathcal{D}$ vary~\cite{vardi1982complexity}.  Table~\ref{tab:complexity} summarises the
results; each row adds one feature to the row above, so the bounds can only grow from one row
to the next.

\begin{table}[!ht]
\centering
\caption{Complexity of EQL query evaluation by feature class.}
\label{tab:complexity}
\renewcommand{\arraystretch}{1.3}
\begin{tabular}{@{} p{6.2cm} p{3.0cm} p{4.6cm} @{}}
\toprule
\textbf{Feature class} & \textbf{Data complexity} & \textbf{Combined complexity} \\
\midrule
CQ and UCQ & $\mathrm{AC}^0$ & NP-complete \\
$+$ NAF and $\forall$ (first-order logic, active domain) & $\mathrm{AC}^0$ & PSPACE-complete \\
$+$ aggregation & $\mathrm{TC}^0$ & PSPACE-complete$^{a}$ \\
$+$ predicates and symbolic functions (as oracles) & $\mathrm{FP}^{\mathcal{O}}$ & PSPACE-complete$^{b}$ \\
\bottomrule
\noalign{\vspace{4pt}}
\multicolumn{3}{@{}l@{}}{%
  \parbox{\dimexpr\linewidth-2\tabcolsep}{%
    \footnotesize $^a$\,For \texttt{count}, \texttt{sum}, \texttt{min}, \texttt{max} and
    \texttt{avg} over numbers of polynomial size.\\[2pt]
    $^b$\,If every callable runs in polynomial time and returns a result of polynomial size.}}
\end{tabular}
\end{table}

\noindent The entries follow from established results:

\begin{itemize}
  \item \textbf{CQ and UCQ.}  CQ evaluation is NP-complete in combined
        complexity~\cite{chandra1977optimal}; UCQ shares the same bounds, and both are in
        $\mathrm{AC}^0$ in data complexity~\cite{abiteboul1995foundations}.

  \item \textbf{$+$ NAF and $\forall$.}  With negation and universal quantification over
        finite domains, query bodies are first-order formulas evaluated under the
        active-domain semantics.  Their combined complexity is
        PSPACE-complete~\cite{vardi1982complexity}: nested alternations of $\forall$ and
        $\exists$ encode quantified Boolean formulas, so unfolding $\forall x\in D$ into a
        finite conjunction does not lower the bound.  Their data complexity is in
        $\mathrm{AC}^0$~\cite{immerman1987languages,barrington1990uniformity}.

  \item \textbf{$+$ aggregation.}  Counting aggregates can express majority, which is not in
        $\mathrm{AC}^0$; first-order logic with aggregates has data complexity in
        $\mathrm{TC}^0$~\cite{hella2001aggregate}.

  \item \textbf{$+$ predicates and symbolic functions.}  We treat callables as oracles.  A
        query $\phi$ with $k$ variables, including quantified ones, whose domains have size at
        most $d$, invokes callables at most $|\phi|\cdot d^{k}$ times.  For a fixed query this
        is polynomial in $|\mathcal{D}|$, so evaluation is in $\mathrm{FP}^{\mathcal{O}}$, and
        it runs in polynomial time whenever every callable does and returns a result of
        polynomial size.  Termination of arbitrary callables cannot be guaranteed
        (Section~\ref{sec:eql:vocab}).
\end{itemize}

\noindent In practice, most EQL queries involve at least one predicate or symbolic function,
so the last row is the operative bound.

\clearpage

% ============================================================
\begin{thebibliography}{99}
% ============================================================

\bibitem{baader2003description}
F.~Baader, D.~Calvanese, D.~McGuinness, D.~Nardi, and P.~Patel-Schneider, editors.
\newblock \emph{The Description Logic Handbook: Theory, Implementation, and
  Applications}.
\newblock Cambridge University Press, 2003.

\bibitem{abiteboul1995foundations}
S.~Abiteboul, R.~Hull, and V.~Vianu.
\newblock \emph{Foundations of Databases}.
\newblock Addison-Wesley, 1995.

\bibitem{apt1988declarative}
K.~R.~Apt, H.~A.~Blair, and A.~Walker.
\newblock Towards a theory of declarative knowledge.
\newblock In J.~Minker, editor, \emph{Foundations of Deductive Databases and
  Logic Programming}, pages 89--148. Morgan Kaufmann, 1988.

\bibitem{chandra1977optimal}
A.~K.~Chandra and P.~M.~Merlin.
\newblock Optimal implementation of conjunctive queries in relational databases.
\newblock In \emph{Proc.\ 9th Annual ACM Symposium on Theory of Computing
  (STOC)}, pages 77--90. ACM, 1977.

\bibitem{clark1978negation}
K.~L.~Clark.
\newblock Negation as failure.
\newblock In H.~Gallaire and J.~Minker, editors, \emph{Logic and Data Bases},
  pages 293--322. Plenum Press, 1978.

\bibitem{kowalski1971linear}
R.~Kowalski and D.~Kuehner.
\newblock Linear resolution with selection function.
\newblock \emph{Artificial Intelligence}, 2(3--4):227--260, 1971.

\bibitem{klug1982equivalence}
A.~C.~Klug.
\newblock Equivalence of relational algebra and relational calculus query
  languages having aggregate functions.
\newblock \emph{Journal of the ACM}, 29(3):699--717, 1982.

\bibitem{lloyd1987foundations}
J.~W.~Lloyd.
\newblock \emph{Foundations of Logic Programming}.
\newblock Springer-Verlag, 2nd edition, 1987.

\bibitem{reiter1978closed}
R.~Reiter.
\newblock On closed world data bases.
\newblock In H.~Gallaire and J.~Minker, editors, \emph{Logic and Data Bases},
  pages 55--76. Plenum Press, 1978.

\bibitem{immerman1987languages}
N.~Immerman.
\newblock Languages that capture complexity classes.
\newblock \emph{SIAM Journal on Computing}, 16(4):760--778, 1987.

\bibitem{barrington1990uniformity}
D.~A.~Mix Barrington, N.~Immerman, and H.~Straubing.
\newblock On uniformity within {NC}$^1$.
\newblock \emph{Journal of Computer and System Sciences}, 41(3):274--306, 1990.

\bibitem{hella2001aggregate}
L.~Hella, L.~Libkin, J.~Nurmonen, and L.~Wong.
\newblock Logics with aggregate operators.
\newblock \emph{Journal of the ACM}, 48(4):880--907, 2001.

\bibitem{rice1953classes}
H.~G.~Rice.
\newblock Classes of recursively enumerable sets and their decision problems.
\newblock \emph{Transactions of the American Mathematical Society},
  74(2):358--366, 1953.

\bibitem{whitehead1910principia}
A.~N.~Whitehead and B.~Russell.
\newblock \emph{Principia Mathematica}, volume~1.
\newblock Cambridge University Press, 1910.

\bibitem{vardi1982complexity}
M.~Y.~Vardi.
\newblock The complexity of relational query languages.
\newblock In \emph{Proc.\ 14th Annual ACM Symposium on Theory of Computing
  (STOC)}, pages 137--146. ACM, 1982.

\end{thebibliography}

\end{document}
